%%%%%%%%%%%%%%%%%%%%%%%%%%%%%%%%%%%%%%%%%
%% Conclusion section
\section{Summary and Concluding Remarks}
\label{sec:conclusion}
This paper uses quasi-random variation in genetic inheritance to estimate the causal effects of the Ed PGI on education and labour market outcomes, and decomposes those effects into a direct genetic channel and an indirect channel operating through education years in a causal mediation framework.
The Ed PGI raises both years of education and later-life income by multiple measures, replicating the core findings of the recent literature \citep{carvalho2024genetics}.
A causal mediation analysis uses these findings as a basis, and at plausible returns to education from the economics literature for Britain, the majority of the total Ed PGI income effect operates through the education years channel, with a residual direct genetic effect that is small and, under higher return values, indistinguishable from zero.
The literature had noted suggestive evidence that direct effects may exist, but had not delivered a quantitative decomposition; this paper does so, and finds that the education years channel accounts for the bulk of the total effect across the plausible range of returns, with (at most) small direct effects.
% The institutions governing access to education are therefore the main mechanism through which genetic endowment shapes earnings inequality, and the main policy lever through which that inequality could be addressed.
